% ============================================================================
% Dok. 267 -- Kosmologische Entartung: ΛCDM und FFGFT als zwei Lesarten
%             derselben Beobachtungen
%
% Zweck: Ausführliche, strukturierte Gegenüberstellung der beiden Modelle.
%        Kernthese: Die klassischen kosmologischen Messungen (SNe Ia, BAO,
%        CMB-Temperatur) sind entartet -- sie können nicht zwischen
%        metrischer Expansion (ΛCDM) und fraktaler Zeitentfaltung (FFGFT)
%        unterscheiden. Die Entscheidung ist theoretisch, nicht empirisch.
%
% Status: Konzeptuelle Klarstellung. Baut auf P18-P20 (Dok. 190) auf.
%         Keine neue Physik, sondern saubere Einordnung des Verhältnisses
%         der beiden Modelle.
% ============================================================================

\section*{Zweck und Kernthese}

Dieses Dokument stellt die beiden kosmologischen Modelle -- das $\Lambda$CDM-
Standardmodell und FFGFT -- in ihrer jeweiligen inneren Logik gegenüber.
Die Kernthese lautet:

\begin{quote}
Die wichtigsten kosmologischen Standardmessungen (Supernovae Typ Ia,
baryonische akustische Oszillationen, CMB-Temperatur) sind \emph{entartet}:
Sie können prinzipiell nicht zwischen einer physikalischen Bewegung des
Raumes ($\Lambda$CDM) und einer intrinsischen, nicht-linearen Entfaltung
der Zeitmetrik (FFGFT) unterscheiden -- solange beide Modelle dieselbe
Abstands-Rotverschiebungs-Beziehung $z(d)$ erzeugen.
\end{quote}

Diese Aussage ist keine Behauptung der Überlegenheit von FFGFT. Sie ist
eine methodische Feststellung: Die Daten allein entscheiden nicht zwischen
den Modellen. Die Entscheidung verlagert sich auf die Ebene der
theoretischen Konsistenz, der Parameteranzahl und der Ableitbarkeit aus
ersten Prinzipien.

% ============================================================================
\section{Die zwei Modelle in ihrer Grundstruktur}
% ============================================================================

\subsection*{Die invertierte Anfangsbedingung}

Der fundamentale Unterschied liegt nicht in den Vorhersagen, sondern in
der Richtung der Entwicklung:

\begin{center}
\begin{tabular}{p{2.3cm}p{5.3cm}p{5.3cm}}
\toprule
\textbf{Phase} & \textbf{$\Lambda$CDM (Standardmodell)} & \textbf{FFGFT} \\
\midrule
Anfang &
  Singularität, Urknall, extrem heiß und dicht, Zeit vergeht schnell &
  Minimale Länge $L_0 = \xi \cdot \ell_P$, fast kompaktifiziert,
  konzentriert aber kalt, Zeit vergeht extrem langsam \\
\addlinespace
Entwicklung &
  Metrische Expansion des Raumes + Abkühlung &
  Fraktale Entfaltung der Zeit $\to$ scheinbare Expansion \\
\addlinespace
Heute &
  Expandiert, kalt, beschleunigt durch dunkle Energie ($\Lambda$) &
  Scheinbar expandiert, stationärer Quanten-Grundzustand, scheinbar
  beschleunigt durch nicht-lineare Zeitrekursion \\
\bottomrule
\end{tabular}
\end{center}

Die beiden Richtungen sind vollständig invertiert:
\begin{itemize}
  \item $\Lambda$CDM: \emph{heiß und schnell} $\to$ \emph{kalt und beschleunigt}
  \item FFGFT: \emph{kalt und langsam} $\to$ \emph{scheinbar warm und beschleunigt}
        (durch Zeitentfaltung)
\end{itemize}

\subsection*{Die Temperatur in beiden Bildern}

In $\Lambda$CDM ist die heutige CMB-Temperatur ($\approx 2{,}7$\,K) der
abgekühlte Überrest einer heißen Frühphase. In FFGFT ist sie kein
Überbleibsel, sondern ein stationärer Gleichgewichtswert aus der
$\xi$-Feld-Quantisierung (Dok.~025): die Strahlung kühlt nicht ab, sie
ist ein inhärentes Merkmal des Vakuums im Gleichgewicht.

In FFGFT ist die anfängliche Konzentration \emph{topologisch}
(Windungszahlen nahe beieinander), nicht \emph{thermodynamisch} (heiß).
Es gibt keine kinetische Energie in großem Maßstab am Anfang, daher keine
hohe Temperatur. Das ist die konsequente Umkehrung der $\Lambda$CDM-
Anfangsbedingung.

% ============================================================================
\section{Die kosmologische Entartung -- drei Beispiele}
% ============================================================================

\subsection{Supernovae Typ Ia und die Lichtkurven-Dehnung}

\textbf{Das Signal:} Die Messdaten zeigen $b \approx 1$: eine ferne
Supernova läuft für uns langsamer ab (Lichtkurve gedehnt), und ihr Licht
ist rotverschoben.

\textbf{$\Lambda$CDM:} Der Raum dehnt sich während des Photonenflugs
metrisch aus. Das streckt die Wellenlänge (Rotverschiebung) und zieht die
Signaldauer auseinander (Zeitdilatation).

\textbf{FFGFT:} Keine Raumexpansion. Aber die fraktale Zeitentfaltung
verlängert den Pfad, den das Signal durch die Geometrie zurücklegt. Diese
geometrische Wegverlängerung dehnt Wellenlänge und Pulsdauer im exakt
selben Verhältnis -- daher ebenfalls $b = 1$.

\textbf{Die Ununterscheidbarkeit:} Beide Mechanismen liefern exakt dieselbe
mathematische Kurve. Die Supernova ,,weiß`` nicht, ob der Raum gewachsen
ist oder ob sich die Zeit fraktal entfaltet hat. Der DES-Wert
$b = 1{,}003 \pm 0{,}011$ schließt ein \emph{nicht-dilatierendes} statisches
Modell ($b = 0$) bei $\sim 90\sigma$ aus -- aber FFGFT sagt $b = 1$ voraus,
nicht $b = 0$. FFGFT ist daher von diesem Test nicht betroffen.

\subsection{Baryonische akustische Oszillationen und die Winkelskalen}

\textbf{Das Signal:} Man misst den Winkelabstand $\theta(z)$ von
Galaxienclustern.

\textbf{$\Lambda$CDM:} Der scheinbare Winkel ändert sich, weil der
Skalenfaktor $a(t)$ wächst und die dunkle Energie diese Expansion
beschleunigt.

\textbf{FFGFT:} Die scheinbare Beschleunigung ist das Resultat der
nicht-linearen zeitlichen Rekursion. Da sich die effektive Zeitrate pro
Distanzschritt ändert, krümmt sich die $z(d)$-Relation genau so, als gäbe
es eine beschleunigte Expansion.

\textbf{Die Ununterscheidbarkeit:} Solange die FFGFT-Zeitentwicklung
dieselbe $z(d)$-Beziehung erzeugt, liefert sie dieselben Winkel am Himmel.
Die BAO-Daten enthalten keinen Marker, der sagt, welcher Mechanismus die
Kurve erzeugt hat.

\textbf{Wichtige Einschränkung (P20):} Der BAO-Test verwendet $D_M/r_s$,
wobei $r_s \approx 147$\,Mpc unter $\Lambda$CDM-Annahmen berechnet wurde.
Onurs $\Delta\chi^2 \approx 138$ vergleicht eine logarithmische
Distanzformel gegen $\Lambda$CDM-Pipeline-Daten -- also Modell gegen Modell,
nicht Theorie gegen rohe Natur.

\subsection{Der kosmische Mikrowellenhintergrund und die Akustikpeaks}

\textbf{Das Signal:} Ein nahezu perfektes Schwarzkörperspektrum bei
$\approx 2{,}7$\,K mit winzigen Temperaturschwankungen und einem
charakteristischen Winkelleistungsspektrum mit Akustikpeaks bei
$\ell \approx 200$, $500$, $800$, \ldots

\subsubsection*{ΛCDM: Hydrodynamische Schallwellen im Urplasma}

In den ersten $\approx 380{.}000$ Jahren nach dem Urknall bildeten
Elektronen, Protonen und Photonen ein vollständig ionisiertes,
undurchsichtiges Photon-Baryon-Fluid. In diesem Medium breiteten sich
Schallwellen aus, angetrieben durch den Wettbewerb zweier Kräfte:

\begin{itemize}
  \item \textbf{Gravitation (Kompression):} Regionen mit geringfügig
        höherer Dichte ziehen Materie an, die Dichte steigt.
  \item \textbf{Strahlungsdruck (Rarefaktion):} Je dichter das Fluid
        komprimiert wird, desto stärker steigt der thermische
        Strahlungsdruck, der schließlich die Gravitation überwindet und
        das Plasma nach außen treibt.
\end{itemize}

Bei der \emph{Rekombination} ($T \approx 3000$\,K) verbanden sich
Elektronen und Protonen zu neutralem Wasserstoff. Die Photonen
entkoppelten schlagartig -- die Schallwellen wurden eingefroren, ihr
Muster als Temperaturschwankung im CMB fixiert.

\textbf{Die Peaks:}
\begin{itemize}
  \item \textbf{1. Peak ($\approx 1°$):} Regionen, die seit dem Urknall
        exakt einmal maximale Kompression erreicht hatten.
  \item \textbf{2. Peak:} Regionen im Zustand maximaler Rarefaktion
        (ein halber Zyklus).
  \item \textbf{3. Peak:} Regionen bei zweiter maximaler Kompression.
\end{itemize}

Die Positionen und Amplituden der Peaks hängen hochempfindlich von den
kosmologischen Parametern ab: das Baryon-zu-Photon-Verhältnis, die
Dunkle-Materie-Dichte, und die Geometrie des Raumes (flach bei $\approx 1°$
für den ersten Peak). $\Lambda$CDM erklärt diesen Datensatz mit
außerordentlicher mathematischer Präzision.

\subsubsection*{FFGFT-Ansatz: Diskrete Gitterresonanzen (strukturell plausibel,
quantitativ offen)}

FFGFT hat keine heiße Frühphase, kein Photon-Baryon-Plasma, keine
Rekombination. Die CMB-Temperatur ist der stationäre Quanten-Grundzustand
des $\Phi$-Feldes (Dok.~025). Die Akustikpeaks müssten durch diskrete
T$^4$-Gitterresonanzen entstehen -- analog zu Emissionsspektren eines
Atoms oder Phonon-Dispersionskurven in einem Kristall.

Aus der geometrischen Resonanzbedingung
\[
  (1+z_*)^{1/\ln(1/\xi)} = \frac{\lambda_e}{L_0},
  \qquad \ln(1/\xi) = 8{,}92,
\]
mit der reduzierten Compton-Wellenlänge $\lambda_e$ und $L_0 = \xi\cdot\ell_P$
folgt eine charakteristische Skala $z_* \approx 875$. Dieser Wert ist eine
rein geometrische FFGFT-Größe (keine Rekombination), die in der Größenordnung
nahe an der $\Lambda$CDM-Rekombination ($z \approx 1089$) liegt, ohne mit ihr
identisch zu sein.

\begin{center}
\begin{tabular}{p{2.3cm}p{5.2cm}p{5.2cm}}
\toprule
\textbf{Aspekt} & \textbf{$\Lambda$CDM} & \textbf{FFGFT-Ansatz} \\
\midrule
Peak-Ursprung & Plasma-Schallwelle, eingefroren & T$^4$-Gitterresonanz \\
Peak-Positionen & Schallhorizont $r_s$ & Torus-Eigenmoden $r/L_\mathrm{tor}$ \\
Peak-Amplituden & Baryon/Photon-Verhältnis & Bose-Einstein-Besetzung \\
\bottomrule
\end{tabular}
\end{center}


% ============================================================================
\section{Die Peak-Struktur aus reiner Geometrie}
% ============================================================================

Es stellt sich die grundsätzliche Frage: Welche Peak-Struktur erzeugt die
reine T$^4$-Geometrie -- ohne Medium, ohne Plasma, ohne den Versuch,
$\Lambda$CDM-Werte zu treffen?

Akustische Resonatoren sind exakt berechenbar, und ihre Eigenmoden-Verhältnisse
hängen \emph{nur von der Geometrie} des Resonators ab, nicht vom Medium:
\begin{center}
\begin{tabular}{p{5.5cm}p{6cm}}
\toprule
\textbf{Resonator} & \textbf{Moden-Verhältnisse} \\
\midrule
Rohr (beide Enden gleich) / Kugel ($j_0$) & $1:2:3:4:5$ (harmonisch) \\
Halboffenes Rohr & $1:3:5:7:9$ (ungerade) \\
Kreismembran ($J_0$-Nullstellen) & $1:2{,}30:3{,}60:4{,}90$ \\
T$^4$-Wicklungsmoden $r=\sqrt{\text{ganz}}$ & $1:1{,}41:1{,}73:2:\ldots$ (dicht) \\
\midrule
\textbf{CMB beobachtet} & $1:2{,}44:3{,}68:5{,}09:6{,}56$ \\
\bottomrule
\end{tabular}
\end{center}

Die Peak-Struktur ist durch die \emph{Geometrie} des Resonanzraums bestimmt,
nicht durch das Medium. $\Lambda$CDM: Resonanzraum = Schallhorizont im Plasma.
FFGFT: Resonanzraum = T$^4$-Torusgeometrie. Beide können dieselbe Serie
erzeugen, wenn die effektive Resonatorgeometrie gleich ist.

\textbf{Diskrete Peaks erfordern eine kompakte Topologie.} Ein kompakter
(endlicher) Resonator erzeugt ein diskretes Modenspektrum; ein
unendlich-offener Raum hätte ein kontinuierliches -- keine scharfen Peaks.
Die beobachteten diskreten CMB-Peaks sprechen daher \emph{für} den kompakten
T$^4$-Torus und \emph{gegen} ein fraktal-unendlich-offenes Universum. Über
die Innen/Außen-Asymmetrie (P20): von außen kompakt-endlich (diskrete Moden),
von innen fraktal-granuliert bis zur Untergrenze $L_0 = \xi\cdot\ell_P$.

% ============================================================================
\section{Der Resonator-Mechanismus: kompaktes Rohr mit Randanpassung}
% ============================================================================

Die reine T$^4$-Wicklungsmoden-Struktur ist zu dicht; eine 1D-stehende Welle
gibt $1:2:3$. Die beobachtete Serie liegt dazwischen und verlangt einen
Versatz. Wie bei einem Blechblasinstrument, wo der Schalltrichter (eine
randbedingungsabhängige Mündungskorrektur) die Stimmung von $1:2:3$
verschiebt: Das kompakte T$^4$ ist das ,,Rohr``, der Übergang zur fraktalen
Innenentfaltung der ,,Schalltrichter``. Das ergibt:
\[
  \boxed{\;\ell_n = \frac{A\,n}{1 + B/n}\;}
\]
mit $A$ der Grundresonanzskala (dem \emph{einen externen Parameter}, der
Torusgröße $= R_H$; aus den Daten $A = 303{,}6 \pm 6{,}8$) und $B$ dem
Versatz. Die \emph{Form} folgt aus ,,kompakter Resonator + randabhängige
Korrektur``; die Größenordnung $B \sim 1/(D-1) = 1/3$ folgt aus der Dimension
$D=4$. $B = 1/3$ ist \emph{keine} externe Eingabe -- es ist eine reine
Dimensionszahl. Nur $A$ ist der externe Parameter.

\begin{center}
\begin{tabular}{p{2.5cm}p{4.5cm}p{4.5cm}}
\toprule
\textbf{Größe} & \textbf{Was sie ist} & \textbf{Externe Eingabe?} \\
\midrule
$A$ (Grundskala) & Torusgröße ($=R_H$) & \textbf{Ja} -- der eine externe
  Parameter (P20) \\
\addlinespace
$B$ (Versatz) & $1/(D-1)=1/3$ aus $D=4$ & \textbf{Nein} -- reine
  Dimensionszahl \\
\bottomrule
\end{tabular}
\end{center}

\emph{Skripte:} \texttt{ffgft\_trompeten\_mechanismus.py},
\texttt{ffgft\_eingaben\_buchhaltung.py}.

% ============================================================================
\section{Der modellneutrale Test: rohe Peak-Verhältnisse}
% ============================================================================

Da $A$ ein externer Parameter ist, kürzt er sich in den Verhältnissen
$\ell_n/\ell_1$ heraus, die dann nur noch von $B=1/3$ (Geometrie) abhängen --
keine externe Skala, kein $H_0$, kein $r_s$, kein Fit:
\begin{center}
\begin{tabular}{p{3.5cm}cccc}
\toprule
 & $\ell_2/\ell_1$ & $\ell_3/\ell_1$ & $\ell_4/\ell_1$ & $\ell_5/\ell_1$ \\
\midrule
FFGFT ($B=1/3$, aus $D=4$) & $2{,}29$ & $3{,}60$ & $4{,}92$ & $6{,}25$ \\
Roh gemessen (Winkel) & $2{,}44$ & $3{,}68$ & $5{,}09$ & $6{,}56$ \\
\bottomrule
\end{tabular}
\end{center}
Mittlere Abweichung $\sim 3{,}3\%$, ohne Planck-Daten -- nur $D=4$ und die
Resonator-Form. \textbf{Einschränkung:} Stammen die Peak-Positionen aus einem
$\Lambda$CDM-Template-Fit, steckt $\Lambda$CDM bereits darin und kürzt sich in
den Verhältnissen \emph{nicht} heraus. Ein sauberer Test braucht modellfreie
lokale Maxima des $C_\ell$-Spektrums (die $\sim 1$--$2\%$ von den
$\Lambda$CDM-Fit-Werten abweichen).

\emph{Skripte:} \texttt{ffgft\_rohe\_peak\_verhaeltnisse.py},
\texttt{ffgft\_lambda\_im\_messprozess.py}.

% ============================================================================
\section{Umkehrung und die symmetrische Zirkularität}
% ============================================================================

Geht man von den gemessenen Peaks aus, ist der Grundton-Winkel
$\theta_1 = \pi/A \approx 0{,}59°$. Aber $\theta_1 = L_\mathrm{res}/D_A$ ist
ein \emph{Verhältnis}: Eindeutig messbar ist nur das Produkt
$H_0 \cdot L_\mathrm{res} = \theta_1 c/\ln(1+z_*) \approx 458$ km/s, nicht
$H_0$ allein. Das Isolieren von $H_0$ braucht eine extern festgelegte Länge --
den externen Parameter aus P20.

Das macht die Zirkularität \textbf{symmetrisch}:
\begin{center}
\begin{tabular}{p{3.5cm}p{4cm}p{4cm}}
\toprule
 & \textbf{$\Lambda$CDM} & \textbf{FFGFT} \\
\midrule
Gemessen & $\theta_1$ (Winkel) & $\theta_1$ (Winkel) \\
Externe Längenskala & $r_s \approx 147$ Mpc & $R_H$ (Torusgröße) \\
Quelle der Skala & heiße Frühphase (Modell) & $E_H/\hbar$ aus $\xi$ \\
$H_0$ & $67{,}4$ (Output) & $66{,}2$ (gesetzt) \\
Zirkulär? & \textbf{Ja} & \textbf{Ja} \\
\bottomrule
\end{tabular}
\end{center}

Ein gemessener Winkel bestimmt immer nur ein Längenverhältnis. Beide Modelle
müssen eine absolute Skala hineinstecken. Niemand gewinnt aus den CMB-Peaks
ein modellneutrales $H_0$. Der faire Vergleich ist nicht ,,wer ist
zirkelfrei`` (keiner ist es), sondern Sparsamkeit (FFGFT: eine externe Skala;
$\Lambda$CDM: $r_s$ aus 6 Parametern plus BBN-Phase) gegen Verankerung
($\Lambda$CDM besser getestet; FFGFTs Frühphase offen).

\emph{Skripte:} \texttt{ffgft\_zirkularitaet\_symmetrie.py},
\texttt{lcdm\_circularity.py}.


% ============================================================================
\section{Die beschleunigte Dilatation in beiden Modellen}
% ============================================================================

Ein oft übersehener Punkt: Beide Modelle sagen eine \emph{beschleunigte}
Dilatation voraus, nicht eine konstante.

\begin{center}
\begin{tabular}{p{3cm}p{5cm}p{5cm}}
\toprule
\textbf{Aspekt} & \textbf{$\Lambda$CDM} & \textbf{FFGFT} \\
\midrule
Grund & Dunkle Energie ($\Lambda$) & Fraktale zeitliche Entwicklung ($\mathcal{R}$) \\
Mechanismus & Abstoßende Gravitation & Akkumulierte Wegverlängerung \\
$w$ & $\approx -1$ & Kein $w$ (kein $\Lambda$-Feld) \\
Beschleunigung & Real (metrische Expansion) & Scheinbar (emergente Zeitdilatation) \\
\bottomrule
\end{tabular}
\end{center}

In FFGFT wächst die effektive Skalenänderung pro Zeitschritt mit der Anzahl
der Rekursionsschritte: Je weiter wir in die Entfernung (kosmologisch: in
die Vergangenheit) blicken, desto mehr entfaltete Zeit liegt zwischen den
Ereignissen. Die Beobachtung einer beschleunigten Dilatation ist daher kein
Widerspruch zu FFGFT -- sie ist genau das, was FFGFT aus seiner eigenen
Dynamik erwartet. Nur die Erklärung ist eine andere.

Dies erschwert die Unterscheidung zwischen den Modellen zusätzlich: Viele
Beobachtungen (SNe~Ia, BAO) messen primär die Beschleunigung selbst, nicht
ihren Mechanismus.

\subsection*{Die drei Phasen der $\Lambda$CDM-Expansionsgeschichte}

Ein struktureller Unterschied betrifft den \emph{zeitlichen Verlauf} der
Expansionsrate. $\Lambda$CDM verlangt drei aufeinanderfolgende Regime:

\begin{enumerate}
  \item \textbf{Frühphase (Inflation):} eine extrem rasche, exponentiell
        beschleunigte Ausdehnung, getrieben durch ein hypothetisches
        Inflaton-Feld mit negativem Druck, das danach zerfällt.
  \item \textbf{Mittlere Phase (gebremste Expansion):} strahlungs- und
        materiedominiert, $a(t) \sim t^{1/2}$ bzw.\ $t^{2/3}$. Die
        anziehende Gravitation \emph{verlangsamt} die Expansion.
  \item \textbf{Heutige Phase (erneute Beschleunigung):} ab $z \approx 0{,}6$
        dominiert die kosmologische Konstante $\Lambda$ und treibt die
        Expansion wieder beschleunigt an.
\end{enumerate}

Die Abfolge \emph{beschleunigt $\to$ gebremst $\to$ wieder beschleunigt} ist
mathematisch durch die Friedmann-Gleichung erlaubt, erfordert aber
\emph{zwei separate, unabhängig postulierte} Beschleunigungsmechanismen
(Inflaton und $\Lambda$), getrennt durch eine Bremsphase. Hinzu kommt das
Koinzidenzproblem: $\Lambda$ wird gerade in unserer kosmischen Epoche
dominant -- eine Feinabstimmung, die das Modell nicht aus sich heraus
erklärt.

FFGFT kommt demgegenüber mit \emph{einem} Mechanismus aus: die scheinbare
Beschleunigung wächst monoton mit der Rekursionstiefe der fraktalen
Zeitentfaltung; es gibt keine getrennte Inflations- und keine getrennte
$\Lambda$-Epoche. Dies ist ein Vorteil auf der Ebene der theoretischen
Ökonomie, \emph{keine} empirische Aussage -- FFGFT hat für die Frühphase
bislang keine quantitativen Vorhersagen, und ob seine eine
Entfaltungsdynamik die gesamte Expansionsgeschichte korrekt nachbildet, ist
nicht gezeigt.

% ============================================================================
\section{Was die Daten können und was nicht}
% ============================================================================

\begin{center}
\begin{tabular}{p{3cm}p{4.3cm}p{4.3cm}}
\toprule
\textbf{Beobachtung} & \textbf{$\Lambda$CDM} & \textbf{FFGFT} \\
\midrule
SNe~Ia: $b \approx 1$ &
  Expansion dehnt Lichtkurven &
  Geometrische Wegverlängerung dehnt Lichtkurven \\
\addlinespace
BAO: $z(d)$ &
  Skalenfaktor $a(t)$ wächst &
  Akkumulierte fraktale Pfadlänge \\
\addlinespace
Beschleunigung &
  Dunkle Energie ($\Lambda$) &
  Nicht-lineare Zeitrekursion ($\mathcal{R}$) \\
\addlinespace
CMB-Spektrum &
  Abkühlung von heißem Urknall &
  Stationärer Quanten-Grundzustand (Dok.~025) \\
\bottomrule
\end{tabular}
\end{center}

\textbf{Was die Daten nicht können:} Sie können nicht unterscheiden, ob die
Dilatation von metrischer Expansion oder von fraktaler Pfadverlängerung
stammt. Beides ergibt dasselbe beobachtbare Signal.

% ============================================================================
\section{Die Grenze für beide Modelle: keine modellneutrale Skala}
% ============================================================================

Die Entartung wird durch ein zweites Problem verstärkt, das beide Modelle
gleichermaßen betrifft (P20, Dok.~190): Es gibt keine modellneutrale
Bestimmung der kosmologischen Skala $R_H$.

\begin{itemize}
  \item $\Lambda$CDM misst $H_0 \approx 67{,}4$\,km/s/Mpc -- aber dieser
        Wert ist eine Pipeline-Ausgabe, die Expansion, FLRW-Metrik und
        Materiegehalt voraussetzt.
  \item FFGFT leitet $E_H/\hbar \approx 66{,}2$\,km/s/Mpc intern aus
        $\xi^{41/4}$ ab -- aber für die kosmologische Anwendung braucht es
        $R_H$ als externen Parameter, der nicht aus $\xi$ allein folgt.
\end{itemize}

Onurs Tests verwenden $\Lambda$CDM-Pipeline-Größen ($r_s$, $H_0$,
kalibrierte Magnituden) als Referenz. Sie testen daher nicht ,,statisch
gegen expandierend``, sondern ,,$\Lambda$CDM-Pipeline gegen die Abwesenheit
dieser Pipeline``. Das ist kein neutraler Test.

% ============================================================================
\section{Wo die Entscheidung tatsächlich fällt}
% ============================================================================

Wenn die Daten blind für den Mechanismus sind, verlagert sich die
Entscheidung auf die theoretische Ebene:

\begin{enumerate}
  \item \textbf{Parameteranzahl:} FFGFT braucht intern nur das fundamentale
        Verhältnis $\xi = 4/30000$ (plus $R_H$ als externen kosmologischen
        Parameter, P20). $\Lambda$CDM benötigt Parameter für dunkle Energie,
        dunkle Materie, Inflation, Baryonendichte und weitere.
  \item \textbf{Singularitäten:} FFGFT erzeugt keine Singularität (kein
        Urknall, keine unendliche Dichte bei $t=0$). $\Lambda$CDM beginnt
        mit einer Singularität.
  \item \textbf{Erklärung der Zeit:} FFGFT leitet die Entstehung und
        Veränderung des Zeitflusses (die fraktale Entfaltung) aus ersten
        Prinzipien her. $\Lambda$CDM postuliert die Zeit als starre,
        unerklärte Hintergrundbühne.
\end{enumerate}

\textbf{Wichtig -- keine Überlegenheitsbehauptung:} Diese Punkte sind
Vorteile auf der Ebene der theoretischen Ökonomie und Konsistenz. Sie sind
\emph{keine} empirische Widerlegung von $\Lambda$CDM.

% ============================================================================
\section{Die ehrliche wissenschaftliche Position}
% ============================================================================

\begin{itemize}
  \item Es gibt zwei konkurrierende Erklärungen für dieselben Beobachtungen.
  \item Keine ist derzeit empirisch bevorzugt, weil die klassischen
        Messungen (SNe~Ia, BAO, CMB-Temperatur) entartet sind.
  \item Die Daten zwingen uns nicht, uns für ein Modell zu entscheiden.
  \item Die Wahl ist eine Frage der theoretischen Konsistenz und der
        Erklärungskraft jenseits der kosmologischen Standarddaten.
  \item FFGFT ist durch Onurs Tests weder bestätigt noch widerlegt. Die
        Tests zeigen lediglich, dass ein \emph{naives} statisches Modell
        ohne fraktale Zeitentfaltung ($b = 0$, keine Beschleunigung)
        ausgeschlossen ist. FFGFT ist kein solches naives Modell.
\end{itemize}

\textbf{Der Punkt, an dem die Modelle prinzipiell unterscheidbar wären:}
Die Frühphase. Hier ist allerdings Vorsicht geboten -- die Frühphase ist
für \emph{beide} Modelle kaum direkt belegbar, und das wird sie vermutlich
auch bleiben: Niemand kann die ersten Momente des Universums direkt
beobachten; beide Modelle rekonstruieren sie aus heutigen Spuren und vielen
Annahmen. $\Lambda$CDM hatte Jahrzehnte Zeit, für diese Phase ein
detailliertes Erklärungsgebäude zu entwickeln (heißer Urknall, Inflation,
Nukleosynthese) -- erfolgreich in der Anpassung an die Daten, aber selbst
auf zahlreichen Annahmen und freien Parametern aufgebaut. FFGFT müsste die
Lichtelementhäufigkeiten (BBN), die CMB-Akustikpeaks und die Strukturbildung
aus seiner umgekehrten Richtung (kalt, langsam, fraktal entfaltend) erklären
-- das ist eine offene Forschungsfrage, an der bisher wenig gearbeitet ist.

Die Frühphase ist daher zwar im Prinzip die Stelle, an der sich die Modelle
unterscheiden könnten -- aber wegen der eingeschränkten Belegbarkeit für
beide Seiten ist nicht zu erwarten, dass sie zu einer eindeutigen
empirischen Entscheidung führt. Eine quantitative FFGFT-Vorhersage für die
frühen Observablen wäre wünschenswert, steht aber aus.

\noindent\textbf{Bezug zu früheren Dokumenten:} P18--P20 (Dok.~190, Status
der kosmologischen Skala), Dok.~025 (CMB als $\xi$-Feld-Grundzustand),
Dok.~182 (fraktale Wegverlängerung), Dok.~264 (Ferrotoroidizität, P''=0
als statischer Grundzustand), Dok.~266 (Analyse von Onurs Tests).
